%% Template para dissertação/tese na classe UFPEthesis
%% versão 0.9.2
%% (c) 2005 Paulo G. S. Fonseca
%% www.cin.ufpe.br/~paguso/ufpethesis

%% Carrega a classe ufpethesis
%% Opções: * Idiomas
%%           pt   - português (padrão)
%%           en   - inglês
%%         * Tipo do Texto
%%           bsc  - para monografias de graduação
%%           msc  - para dissertações de mestrado (padrão)
%%           qual - exame de qualificação doutorado
%%           prop - proposta de tese doutorado
%%           phd  - para teses de doutorado
%%         * Mídia
%%           scr  - para versão eletrônica (PDF) / consulte o guia do usuario
%%         * Estilo
%%           classic - estilo original à la TAOCP (deprecated)
%%           std     - novo estilo à la CUP (padrão)
%%         * Paginação
%%           oneside - para impressão em face única
%%           twoside - para impressão em frente e verso (padrão)
\documentclass{ufpethesis}

%% Preâmbulo:
%% coloque aqui o seu preâmbulo LaTeX, i.e., declaração de pacotes,
%% (re)definições de macros, medidas, etc.

%% Identificação:

% Universidade
% e.g. \university{Universidade de Campinas}
% Na UFPE, comente a linha a seguir
\university{Universidade Federal de Pernambuco}

% Endereço (cidade)
% e.g. \address{Campinas}
% Na UFPE, comente a linha a seguir
\address{Recife}

% Instituto ou Centro Acadêmico
% e.g. \institute{Centro de Ciências Exatas e da Natureza}
% Comente se não se aplicar
\institute{Centro de Informática}

% Departamento acadêmico
% e.g. \department{Departamento de Informática}
% Comente se não se aplicar
\department{Programa de Pós-graduação em Ciência da Computação}

% Programa de pós-graduação
% e.g. \program{Pós-graduação em Ciência da Computação}
\program{Programa de Pós-graduação em Ciência da Computação}

% Área de titulação
% e.g. \majorfield{Ciência da Computação}
\majorfield{Ciência da Computação}

% Título da dissertação/tese
% e.g. \title{Sobre a conjectura $P=NP$}
\title{Avaliação de Modelos Não Supervisionados via Habilidade Latente e Concordância entre Métodos}

% Data da defesa
% e.g. \date{19 de fevereiro de 2003}
\date{2026}

% Autor
% e.g. \author{José da Silva}
\author{Manuel Ferreira Junior}

\concentrationarea{TEMPPPPPPPPPP}

% Orientador(a)
% Opção: [f] - para orientador do sexo feminino
% e.g. \adviser[f]{Profa. Dra. Maria Santos}
\adviser{Ricardo Bastos Cavalcante Prudêncio}

% Orientador(a)
% Opção: [f] - para orientador do sexo feminino
% e.g. \coadviser{Prof. Dr. Pedro Pedreira}
% Comente se não se aplicar
\coadviser{Telmo de Menezes e Silva Filho}

%% Inicio do documento
\begin{document}

%%
%% Parte pré-textual
%%
\frontmatter

% Folha de rosto
% Comente para ocultar
\frontpage

% Portada (apresentação)
% Comente para ocultar
\presentationpage

% Dedicatória
% Comente para ocultar
\begin{dedicatory}
Este trabalho é dedicado a Deus, por tudo que foi desenvolvido; a minha namorada, Hayse, que sempre esteve ao meu lado desde antes; aos meus pais, Manuel e Rosilene, pelo auxilio durante toda a minha vida; aos demais familiares e amigos, que de forma direta ou indiretamente contribuíram para o meu progresso até este momento.
%\begin{epigraph}[]{}%{Francis Drake}

\vspace{0.5cm}

\textit{ Sic Parvis Magna}\\
\small “A grandeza nasce de pequenos inícios.”
%\end{epigraph}
\end{dedicatory}

% Agradecimentos
% Se preferir, crie um arquivo à parte e o inclua via \include{}
\acknowledgements
<DIGITE OS AGRADECIMENTOS AQUI>

% Epígrafe
% Comente para ocultar
% e.g.
%  \begin{epigraph}[Tarde, 1919]{Olavo Bilac}
%  Última flor do Lácio, inculta e bela,\\
%  És, a um tempo, esplendor e sepultura;\\
%  Ouro nativo, que, na ganga impura,\\
%  A bruta mina entre os cascalhos vela.
%  \end{epigraph}
\begin{epigraph}[]{Immanuel Kant}
\textit{Sapere aude}\\
\small “Ouse saber.”
\end{epigraph}

% Resumo em Português
% Se preferir, crie um arquivo à parte e o inclua via \include{}
\resumo
% A avaliação de modelos não supervisionados é um problema central em Aprendizado de Máquina, especialmente na ausência do rótulos real associado a instância. Essa dissertação propõe uma abordagem unificada baseada na Teoria de Resposta ao Item (TRI) para estimar habilidade latente de modelos e dificuldade de instâncias, explorando a concordância entre métodos para como se agrupar pares de instâncias. Inicialmente, é apresentado o modelo $\beta^4$-IRT, sendo este uma extensão do $\beta^3$-IRT, decompondo o parâmetro de discriminação em duas partes associadas ao sinal e magnitude, mitigando problemas de não-identificabilidade, melhorando assim a recuperação dos parâmetros, em especial sobre instâncias ruidosas. Em seguida, é proposto o método  CLAIRE, formulando a avaliação de agrupamentos como um problema de estimação de habilidade via TRI, utiliznaod uma matriz de resposta construída a partir da concordância entre modelos quanto ao agruapmento de pares de instâncias. Os resultados experimentais demonstram que a abordagem permite ranquear modelos de clustering sem acesso ao rótulo real, apresentando uma alta correlação com medidas externas tradicionais. Além disso, experimentos adicionais evidenciam que, mesmo com a inserção de partições aleatórias no conjunto de modelos avaliados, o CLAIRE mantém um ranqueamento consistente e robusto, preservando a ordenação correta dos modelos informativos no pool considerado. Assim, a dissertação estabelece a modelagem de habilidade latente como um princípio robusto e geral para avaliação de métodos não supervisionados.

A avaliação de modelos não supervisionados é um problema central em Aprendizado de Máquina, especialmente na ausência de rótulos reais associados às instâncias. Diferentemente do cenário supervisionado, em que métricas externas permitem mensurar diretamente o desempenho dos modelos, métodos não supervisionados dependem, em geral, de medidas internas baseadas em coesão e separação dos agrupamentos, as quais são fortemente influenciadas pela métrica de distância adotada e nem sempre refletem a capacidade do modelo de recuperar estruturas latentes ocultas. Nesse contexto, torna-se necessário o desenvolvimento de abordagens mais gerais e robustas para avaliação de modelos de agrupamento. Esta dissertação propõe uma abordagem unificada baseada na Teoria de Resposta ao Item (TRI) para estimar habilidade latente de modelos e dificuldade de instâncias, explorando a concordância entre métodos quanto ao agrupamento de pares de instâncias como fonte de informação avaliativa. Inicialmente, é apresentado o modelo $\beta^4$-IRT, uma extensão do $\beta^3$-IRT, no qual o parâmetro de discriminação é decomposto em duas componentes associadas ao sinal e à magnitude. Essa reformulação mitiga problemas de não-identificabilidade presentes no modelo original e melhora significativamente a recuperação dos parâmetros, particularmente em cenários com instâncias ruidosas ou com discriminação negativa, comuns em aplicações de Aprendizado de Máquina. Em seguida, é proposto o método CLAIRE, que formula a avaliação de agrupamentos como um problema de estimação de habilidade via TRI, utilizando uma matriz de respostas construída a partir da concordância entre modelos quanto ao agrupamento de pares de instâncias. Ao assumir que modelos mais informativos tendem a concordar entre si sobre decisões estruturais relevantes, o método permite estimar habilidades relativas sem acesso a rótulos verdadeiros. Os resultados experimentais demonstram que a abordagem produz ranqueamentos altamente correlacionados com medidas externas tradicionais. Além disso, experimentos adicionais evidenciam que, mesmo com a inserção de partições aleatórias no conjunto de modelos avaliados, o CLAIRE mantém um ranqueamento consistente e robusto, preservando a ordenação correta dos modelos informativos no pool considerado. Assim, a dissertação estabelece a modelagem de habilidade latente como um princípio teórico e metodológico geral para a avaliação de métodos não supervisionados.
% Palavras-chave do resumo em Português
\begin{keywords}
Agrupamento, Avaliação, Teoria de Resposta ao Item, Concordância entre modelos, Aprendizado de Máquina.
\end{keywords}

% Resumo em Inglês
% Se preferir, crie um arquivo à parte e o inclua via \include{}
\abstract

The evaluation of unsupervised models is a central problem in Machine Learning, particularly in the absence of ground-truth labels associated with instances. Unlike the supervised setting, where external metrics allow direct measurement of model performance, unsupervised methods generally rely on internal measures based on cluster cohesion and separation, which are strongly influenced by the adopted distance metric and do not necessarily reflect the model’s ability to recover hidden latent structures. In this context, the development of more general and robust approaches for clustering evaluation becomes necessary. This dissertation proposes a unified approach based on Item Response Theory (IRT) to estimate latent model ability and instance difficulty, exploring agreement between methods regarding the grouping of instance pairs as a source of evaluative information. Initially, the $\beta^4$-IRT model is introduced as an extension of $\beta^3$-IRT, in which the discrimination parameter is decomposed into two components associated with sign and magnitude. This reformulation mitigates non-identifiability issues present in the original model and significantly improves parameter recovery, particularly in scenarios involving noisy instances or negative discrimination, which are common in Machine Learning applications. Subsequently, the CLAIRE method is proposed, formulating clustering evaluation as an ability estimation problem via IRT, using a response matrix constructed from agreement between models regarding the grouping of instance pairs. By assuming that more informative models tend to agree with each other on structurally relevant decisions, the method enables the estimation of relative abilities without access to ground-truth labels. Experimental results demonstrate that the approach produces rankings highly correlated with traditional external measures. Furthermore, additional experiments show that even with the insertion of random partitions into the pool of evaluated models, CLAIRE maintains a consistent and robust ranking, preserving the correct ordering of informative models within the considered pool. Thus, this dissertation establishes latent ability modeling as a general theoretical and methodological principle for the evaluation of unsupervised methods.
% Palavras-chave do resumo em Inglês
\begin{keywords}
Clustering, Evaluation, Item Response Theory, Model Agreement, Machine Learning.
\end{keywords}

% Sumário
% Comente para ocultar
\tableofcontents

% Lista de figuras
% Comente para ocultar
\listoffigures

% Lista de tabelas
% Comente para ocultar
\listoftables



%%
%% Parte textual
%%
\mainmatter

% É aconselhável criar cada capítulo em um arquivo à parte, digamos
% "capitulo1.tex", "capitulo2.tex", ... "capituloN.tex" e depois
% incluí-los com:
\chapter{Introdução}

A avaliação de modelos é um dos pilares centrais do Aprendizado de Máquina, orientando a seleção de modelos, a interpretação de resultados e o avanço metodlógico da área. No contexto supervisionado, métricas de desempenho podem ser diretamente calculadas a partir de rótulos verdadeiros, uma vez que são conhecidos, permitindo comparações mais objetivas entre desempenho de modelos. Entretanto, no aprendizado não supervisionado, particularmente em tarefas de agrupamento, a ausência de rótulos reais torna o problema de avaliação significativamente mais desafiador. Nesses cenários, a qualidade de uma partição deve ser inferida a partir da própria estrutura dos dados ou da comparação indireta entre modelos, levantando questões sobre robustez, interpretabilidade e generalidade das medidas utilizadas.

Tradicionalmente, a validação de agrupamentos é realizada por meio é realizada por meio de medidas internas ou externas. Medidas externas como o Rand Index ajustado dependem da existência de uma partição de referência e são amplamente utilizadas em experimentos controlados. Contudo, tais medidas tornam-se inviáveis em aplicações reais, onde o objetivo do agrupamento desconhece o rótulo verdadeiros. Medidas internas, por sua vez, avaliam propriedades como coesão e separação dos clusters, mas são sensíveis à métrica de distância adotada, podendo apresentar vieses em relação a determinadas estruturas e frequentemente produzem ranqueamentos contraditórios entre modelos. Além disso, ambas as abordagens negligenciam uma dimensão fundamental: a variabilidade de dificuldade entre instâncias, tratando todos os pontos de dados com uma dificuldade constante.

Nesse contexto, esta dissertação adota a Teoria de Resposta ao Item (TRI) como arcabouço teórico para modelar avaliação no cenário de Aprendizado de Máquina. Originalmente desenvolvida na psicometria para estimar habilidades latentes de respondentes e dificuldades de itens, a TRI permite analisiar desempenho considerando explicitamente a interação entre habilidade e dificuldade. Em anos recentes, modelos baseados em TRI foram adaptados para avaliar sistemas de Inteligência Artificial em tarefas supervisionadas, possibilitando estimar habilidade de modelos e identificar instâncias mais desafiadoras ou podendo ser consideradas como ruidosas. Essa perspectiva oferece uma alternativa às métricas tradicionais, ao introduzir uma modelagem probablística que incorpora um disciminante e uma dificuldade associadas as instâncias, fornecendo interpretações mais refinadas sobre o desempenho.

A primeira contribuição desta dissertação consiste no aprimoramento do uso para o modelo $\beta^{4}$-IRT, uma extensão do $\beta^{3}$-IRT voltada para respostas contínuas limitadas entre 0 e 1. O modelo trata explicitamente o problema de não-identificabilidade associado ao parâmetro de discriminação, fatorando-o em componentes de sinal e magnitude e adotando uma estratégia de estimação que melhora a recuperação de parâmetros. Essa reformulação mostra-se particularmente relevante em cenários de Aprendizado de Máquina, nos quais discriminações negativas podem aparecer em razão de instãncias ruidosas ou ambíguas.

A segunda contribuição é o método CLAIRE (CLuster Agreement-based Item REsponses), que reformula a avaliação de agrupamentos como um problema de estimação de habilidade latente via TRI. Em vez de depender do rótulo real da instância, o método constrói uma matriz de resposta baseada na concordância entre modelos quanto ao agrupamento de pares de instâncias. Sob essa perspectiva, modelos mais informativos são aqueles que tendem a concordar entre si nas instâncias mais difíceis, enquanto instãncias difíceis são aquelas nas quais apenas modelos de alta habilidade apresentam um consenso. Experimentos demosntram que essa abordagem produz ranqueamentos altamente correlacionados com medidas externas tradicionais e permanece robusta mesmo quando partições aleatórias são adicionadas ao conjunto de modelos avaliados.

Ao integrar avanços metodológicos na modelagem de discriminação com uma nova fomulação para validação de agrupamentos, esta dissertação estabelece a habilidade latente como princípio unificador para avaliação de métodos não supervisionados. Dessa forma, contribui tanto par ao desenvolvimento teórico de métricas de avaliação quanto para aplicações práticas em cenários onde a ausência de rótulos é a regra, e não a exceção.
% \include{capitulo2}
% ...
% \include{capituloN}



%%
%% Parte pós-textual
%%
\backmatter

% Apêndices
% Comente se não houver apêndices
\appendix

% É aconselhável criar cada apêndice em um arquivo à parte, digamos
% "apendice1.tex", "apendice.tex", ... "apendiceM.tex" e depois
% incluí-los com:
% \include{apendice1}
% \include{apendice2}
% ...
% \include{apendiceM}


% Bibliografia
% É aconselhável utilizar o BibTeX a partir de um arquivo, digamos "biblio.bib".
% Para ajuda na criação do arquivo .bib e utilização do BibTeX, recorra ao
% BibTeXpress em www.cin.ufpe.br/~paguso/bibtexpress
\nocite{*}
\bibliographystyle{alpha}
\bibliography{sample}

% Cólofon
% Inclui uma pequena nota com referência à UFPEThesis
% Comente para omitir
\colophon

%% Fim do documento
\end{document}
